
\subsubsection{Mechanistic Diagnosis: Batch Contamination}

The inversion effect requires a mechanistic explanation. We propose \textbf{batch contamination} as the primary candidate: under standard ERM, minority samples incur high loss and large gradient norms. With batch size $B=32$ and ~5\% minority prevalence, 1--2 minority samples per batch carry disproportionately large gradient magnitudes. These pull the within-batch mean gradient toward the minority gradient direction, increasing minority samples' measured cosine similarity with the mean --- producing the observed inversion.

Formally, if minority gradients have magnitude $k > 1$ times majority gradient magnitude, the minority contribution to the batch mean is $\propto k \cdot n_\text{min} / B$. With $k$ potentially large at epoch 1 due to high minority loss (empirically, loss ROC-AUC = 0.930 on Waterbirds at epoch 1 indicates minority samples incur substantially higher loss than majority, consistent with large $k$, though $k$ is not directly measured here) and $n_\text{min} = 1$, a single minority sample's gradient contribution to the mean may be comparable to the combined majority contribution.

\textbf{Prevalence evidence:} CelebA's weaker inversion (0.246 vs. 0.150) at much lower minority prevalence (~0.8\%) is quantitatively consistent with contamination being proportional to $k \cdot n_\text{min} / B$.

\subsubsection{Alternative Explanations}

\textbf{Gradient compression at last layer:} The Linear(2048, 2) projects to 2 class logits; group-discriminative signal in earlier 2048-dim representations may be destroyed. Penultimate-layer alignment experiments would distinguish this from batch contamination --- this is an immediate next step.

\textbf{ImageNet pretraining:} At epoch 1, both majority and minority samples may produce similar fine-tuning gradients, reducing directional discriminability. This may compound with contamination at early epochs.

These alternatives are not ruled out by the current experiments. The batch contamination hypothesis is supported by the prevalence-dependent pattern across datasets, but definitive separation from these alternatives requires additional ablations.

\subsubsection{The Expected Corrective Direction}

A \textbf{two-pass global mean gradient} is expected to address the root cause:
\begin{itemize}
\item Pass 1: Standard ERM training
\item $Pass 2: Compute \bar{g}_\text{global} = \frac{1}{N}\sum_{i=1}^N g_i over full training set$
\item Align each sample's gradient against $\bar{g}_\text{global}$ (not per-batch mean)
\end{itemize}

The global mean is stable, majority-dominated across thousands of samples, and robust to per-batch minority gradient spikes. The annotation-free property is preserved. Computational cost: one additional full-dataset pass per epoch. Whether this correction resolves the inversion remains to be empirically validated.

\subsubsection{Limitations}

\begin{itemize}
\item \textbf{Single seed (42):} Appropriate for a proof-of-concept existence test; the failure margin (0.27--0.78 gap) is replication-grade.
\item \textbf{CelebA 16K subsample:} Group composition ratios are preserved by stratified sampling; the conclusion is unchanged.
\item \textbf{Last layer only:} Cannot conclude alignment fails at all layers; penultimate-layer alignment is untested.
\item \textbf{Global mean gradient untested:} Identified as the expected corrective direction, but not yet empirically validated.
\item \textbf{Batch contamination not directly isolated:} The prevalence-dependent pattern across datasets is consistent with contamination, but does not rule out last-layer gradient compression or ImageNet pretraining as alternative or contributing explanations.
\end{itemize}

\subsubsection{Broader Implications}

The batch contamination failure mode may affect any method using within-batch gradient statistics for minority identification in imbalanced settings. Methods should prefer globally computed gradient reference directions over per-batch statistics when minority prevalence is low and minority samples have high loss.

